% Shared style for integrated experiment handouts.
\usepackage[left=2.15cm,right=2.15cm,top=2.25cm,bottom=2.15cm,headsep=0.7cm]{geometry}
\usepackage{amsmath,amssymb}
\usepackage{booktabs,array,tabularx,longtable,multirow}
\usepackage{enumitem}
\usepackage{xcolor}
\usepackage{graphicx}
\usepackage{tikz}
\usetikzlibrary{arrows.meta,positioning,calc,shapes.geometric}
\usepackage{tcolorbox}
\usepackage{listings}
\usepackage{hyperref}
\usepackage{fancyhdr}
\usepackage{chngcntr}

\definecolor{Ink}{HTML}{17233C}
\definecolor{Navy}{HTML}{173B68}
\definecolor{Blue}{HTML}{2878B5}
\definecolor{Cyan}{HTML}{2A9D8F}
\definecolor{Amber}{HTML}{E9A23B}
\definecolor{Red}{HTML}{C84B4B}
\definecolor{Paper}{HTML}{F6F8FB}
\definecolor{PaleBlue}{HTML}{EAF2FA}
\definecolor{PaleCyan}{HTML}{EAF7F4}
\definecolor{PaleAmber}{HTML}{FFF5E4}
\definecolor{Rule}{HTML}{D6DEE8}
\definecolor{Muted}{HTML}{5B677A}

\hypersetup{
  colorlinks=true,
  linkcolor=Navy,
  urlcolor=Blue,
  citecolor=Cyan,
  pdfauthor={《人工智能数学原理与算法》课程组},
  pdftitle={\DocTitle},
  pdfsubject={\DocSubject}
}

\ctexset{
  section={format=\Large\bfseries\color{Navy},beforeskip=1.4ex plus .2ex,afterskip=.8ex},
  subsection={format=\large\bfseries\color{Ink},beforeskip=1.2ex plus .2ex,afterskip=.55ex},
  subsubsection={format=\normalsize\bfseries\color{Blue},beforeskip=1ex,afterskip=.35ex}
}

\setlength{\parindent}{2em}
\setlength{\parskip}{0.22em}
\setlength{\emergencystretch}{2em}
\linespread{1.14}
\setlist[itemize]{leftmargin=2em,itemsep=.25em,topsep=.35em}
\setlist[enumerate]{leftmargin=2.2em,itemsep=.28em,topsep=.35em}
\renewcommand{\arraystretch}{1.30}
\newcolumntype{Y}{>{\raggedright\arraybackslash}X}
\newcolumntype{C}[1]{>{\centering\arraybackslash}p{#1}}
\newcolumntype{L}[1]{>{\raggedright\arraybackslash}p{#1}}

\newcommand{\R}{\mathbb{R}}
\newcommand{\E}{\mathbb{E}}
\newcommand{\MSE}{\operatorname{MSE}}
\newcommand{\RMSE}{\operatorname{RMSE}}
\newcommand{\argmin}{\operatorname*{arg\,min}}
\newcommand{\vect}[1]{\boldsymbol{#1}}
\newcommand{\mat}[1]{\mathbf{#1}}
\newcommand{\term}[1]{\textcolor{Blue}{\textbf{#1}}}
\newcommand{\code}[1]{\texttt{#1}}
\numberwithin{equation}{subsection}

\tcbset{boxrule=.7pt,arc=1.6mm,left=2.5mm,right=2.5mm,top=1.8mm,bottom=1.8mm,before skip=7pt,after skip=7pt}
\newtcolorbox{keybox}[1][核心结论]{colback=PaleBlue,colframe=Blue,title={#1},fonttitle=\bfseries,coltitle=Ink}
\newtcolorbox{examplebox}[1][贯穿例题]{colback=PaleCyan,colframe=Cyan,title={#1},fonttitle=\bfseries,coltitle=Ink}
\newtcolorbox{questionbox}[1][检查点]{colback=PaleAmber,colframe=Amber,title={#1},fonttitle=\bfseries,coltitle=Ink}
\newtcolorbox{pitfallbox}[1][易错提醒]{colback=red!3,colframe=Red,title={#1},fonttitle=\bfseries,coltitle=Ink}
\newtcolorbox{teacherbox}[1][教师提示]{colback=Paper,colframe=Rule,title={#1},fonttitle=\bfseries,coltitle=Muted}
\newtcolorbox{proofbox}[1][推导与证明]{colback=Blue!3,colframe=Navy,title={#1},fonttitle=\bfseries,coltitle=Ink}

\lstdefinestyle{python}{
  language=Python,
  basicstyle=\ttfamily\small,
  keywordstyle=\color{Blue}\bfseries,
  commentstyle=\color{Cyan!70!black},
  stringstyle=\color{Red!80!black},
  showstringspaces=false,
  columns=fullflexible,
  keepspaces=true,
  frame=single,
  rulecolor=\color{Rule},
  backgroundcolor=\color{Paper},
  breaklines=true,
  numbers=left,
  numberstyle=\tiny\color{Muted},
  xleftmargin=1.5em,
  framexleftmargin=1.2em,
  aboveskip=.7em,
  belowskip=.7em
}
\lstset{style=python}

\pagestyle{fancy}
\fancyhf{}
\fancyhead[L]{\small\color{Muted}人工智能数学原理与算法}
\fancyhead[R]{\small\color{Muted}\DocShortTitle}
\fancyfoot[L]{\small\color{Muted}\DocType}
\fancyfoot[R]{\small\color{Muted}\thepage}
\renewcommand{\headrulewidth}{0pt}
\renewcommand{\footrulewidth}{0pt}

\newcommand{\makeexperimenttitle}[4]{
  \hypersetup{pageanchor=false}
  \begin{titlepage}\thispagestyle{empty}
  \begin{tikzpicture}[remember picture,overlay]
    \fill[Ink] (current page.north west) rectangle ([yshift=-9.4cm]current page.north east);
    \fill[Blue] ([yshift=-9.4cm]current page.north west) rectangle ([yshift=-9.68cm]current page.north east);
    \fill[Cyan] ([xshift=2.15cm,yshift=-1.55cm]current page.north west) rectangle ([xshift=2.27cm,yshift=-8.55cm]current page.north west);
    \fill[Paper] ([xshift=2.15cm,yshift=2.0cm]current page.south west) rectangle ([xshift=-2.15cm,yshift=6.15cm]current page.south east);
  \end{tikzpicture}
  \vspace*{1.45cm}
  {\color{white}\Large 人工智能数学原理与算法\par}\vspace{.55cm}
  {\color{white}\fontsize{31}{39}\selectfont\bfseries #1\par}\vspace{.35cm}
  {\color{white}\fontsize{21}{28}\selectfont #2\par}
  \vfill
  \begin{tcolorbox}[colback=white,colframe=white,boxrule=0pt,arc=0mm,left=7mm,right=7mm,top=6mm,bottom=6mm,width=\textwidth]
    \begin{tabularx}{\linewidth}{L{2.6cm}Y}
      \textcolor{Muted}{材料类型} & \DocType\\
      \textcolor{Muted}{实验安排} & #3\\
      \textcolor{Muted}{贯穿案例} & #4\\
      \textcolor{Muted}{适用对象} & 已完成机器学习导论、具备高中数学与基础计算机操作能力的本科生\\
    \end{tabularx}
  \end{tcolorbox}
  \vspace{.7cm}{\color{Muted}\small 版本：2026年秋季学期整合稿\hfill 源文件与 PDF 同目录交付\par}
  \end{titlepage}
  \hypersetup{pageanchor=true}
  \pagenumbering{roman}\tableofcontents\clearpage\pagenumbering{arabic}
}
